\documentclass[leqno]{article}
% Shared mark scheme preamble for resources.danieljsmith.org: the pack
% preamble, unchanged, then the scheme's own layer. A scheme's .tex is
% rendered from its YAML (scripts/README.md, Mark schemes), never edited.
\input{\CurrentFilePath/pack_preamble.tex}

% The page: equal side margins, so the table runs the full width of the
% header rule. The footers' right edge follows the text's.
\geometry{left=0.8in, right=0.8in}
\makeatletter
\renewcommand{\djs@footR}[1]{%
  \put(\dimexpr1in+\hoffset+\oddsidemargin+\textwidth\relax,
    -\dimexpr\paperheight-0.32in\relax){\makebox[0pt][r]{#1}}}
\makeatother

\usepackage[export]{adjustbox}
\usepackage{tabularray}
\definecolor{msGrid}{RGB}{176, 172, 164}
\definecolor{msTint}{RGB}{245, 243, 239}

% The header on every page, the front page included, and the booklet label
% under the contents-page title. \djsQuestionPageStyle stays empty.
\newcommand{\djsSchemeHeader}[1]{%
  \djsPackHeader[markscheme]{\textit{Mark scheme} - #1}%
  \renewcommand{\djsPackTitle}{#1}%
  \renewcommand{\djsBookletLabel}{Mark scheme}%
}

% The front page: the contents heading, spaced as tocloft spaces the packs'
% "Questions", then the contents table. A question's name and page number
% both link to its first page.
\newcommand{\msContentsHeading}{%
  \vskip\cftbeforetoctitleskip
  {\noindent\cfttoctitlefont\contentsname\par}%
  \vskip\cftaftertoctitleskip}
\NewTblrEnviron{mscontents}
\SetTblrInner[mscontents]{
  colspec = {Q[wd=26mm, l] X[l] Q[wd=16mm, c] Q[wd=14mm, r]},
  width = \linewidth,
  rowsep = 4pt, colsep = 6pt,
  hlines = {0.4pt, msGrid}, vlines = {0.4pt, msGrid},
  hline{2} = {2.2pt, scarlet},
  row{1} = {font=\small\bfseries},
  row{Z} = {font=\bfseries},
  column{2} = {font=\small},
}
\newcommand{\msContentsQuestion}[1]{\hyperref[q:#1]{Question #1}}
\newcommand{\msContentsPage}[1]{\hyperref[q:#1]{\pageref*{q:#1}}}
\newcommand{\msSep}{\hspace{0.9em}}

% The conventions pages' layout: headings in the documents' own bold, tables
% ruled and padded as the scheme's, roomy lists.
\newcommand{\msHead}[1]{%
  \par\vspace{18pt}%
  {\noindent\bfseries #1\par}%
  \vspace{6pt}}
\NewTblrEnviron{mstblr}
\SetTblrInner[mstblr]{rowsep = 6pt, colsep = 6pt,
  hlines = {0.4pt, msGrid}, vlines = {0.4pt, msGrid}}
\newenvironment{msList}
  {\begin{itemize}[leftmargin=1.2em, itemsep=5pt, topsep=3pt, parsep=0pt]}
  {\end{itemize}}

% A question's first page: the link target and the PDF bookmark, as
% \questionitem gives the packs.
\newcommand{\msQnum}{}
\newcommand{\msQuestion}[1]{%
  \clearpage
  \renewcommand{\msQnum}{#1}%
  \phantomsection\label{q:#1}%
  \addcontentsline{toc}{subsection}{\protect\djsTocEntry{#1}}}

% The scheme's table: Question 16mm, Scheme the rest, Marks 16mm, Notes
% 44mm; every cell ruled in msGrid; the scarlet rule under the head row. A
% long table has no caption or "continued" foot; a continuation page opens
% with a muted line above the repeated head row. The rules across a part's
% rows, the alternative's tint and the total row are the rendered table's.
\NewTblrTheme{ms}{%
  \DefTblrTemplate{firsthead}{default}{}%
  \DefTblrTemplate{middlehead,lasthead}{default}{%
    {\small\itshape\color{djsMuted}Question \msQnum\ continued\par}%
    \vspace{6pt}}%
  \DefTblrTemplate{firstfoot,middlefoot,lastfoot}{default}{}%
}
\NewTblrEnviron{msscheme}
\SetTblrOuter[msscheme]{long, theme=ms, entry=none, label=none}
\SetTblrInner[msscheme]{
  colspec = {Q[wd=16mm, l] X[l] Q[wd=16mm, c] Q[wd=44mm, l]},
  width = \linewidth,
  rowsep = 6pt, colsep = 6pt,
  rowhead = 1,
  vlines = {0.4pt, msGrid},
  hline{1} = {0.4pt, msGrid}, hline{2} = {2.2pt, scarlet},
  hline{Z} = {0.4pt, msGrid},
  column{1} = {font=\bfseries}, column{3} = {font=\bfseries},
  column{4} = {font=\small},
  row{1} = {rowsep=4pt, font=\small\bfseries},
}

% Two marks on a row: each note starts with its code.
\newcommand{\msNote}[2]{\textbf{#1}\enspace #2}
% A full-row figure: the worked solution's own picture, scaled uniformly and
% only down, to the table less the Question column. The picture is drawn
% with the line width it has in the solutions document, at its list depth
% there (the optional argument, 1 by default: a part's working), so one
% sized by \linewidth keeps its size; tabularray sets the cell's \linewidth
% to the cell's measured width, which would shrink it at every pass. The
% widths are measured once, in a solution box at the packs' text width
% (pack_preamble.tex's margins, 0.8in and 1in).
\newlength{\msSpanW}
\makeatletter
\newcommand{\ms@keepLW}[1]{%
  \expandafter\xdef\csname ms@solnLW@#1\endcsname{\the\linewidth}}
\AtBeginDocument{%
  \setlength{\msSpanW}{\dimexpr\textwidth-16mm-28pt\relax}%
  \setbox\z@\vbox{%
    \begin{minipage}{\dimexpr\paperwidth-1.8in\relax}
      \begin{djsSolution}
        \ms@keepLW{0}%
        \begin{enumerate}\item \ms@keepLW{1}%
        \begin{enumerate}\item \ms@keepLW{2}%
        \begin{enumerate}\item \ms@keepLW{3}%
        \end{enumerate}\end{enumerate}\end{enumerate}
      \end{djsSolution}
    \end{minipage}}}
\newcommand{\msFigure}[2][1]{%
  \adjustbox{max width=\msSpanW, max totalheight=0.6\textheight, valign=t}{%
    \hsize=\csname ms@solnLW@#1\endcsname
    \linewidth=\hsize
    #2}}
\makeatother
% A figure mark's feature, in the Scheme column.
\newcommand{\msFeature}[1]{\textit{#1}}
% An alternative's title row.
\newcommand{\msAlternative}[1]{\textit{Alternative:} #1}
% Drafts only: a queried mark's code, the query's why under its note, and a
% flag in its part's place.
\newcommand{\msQueried}[2]{#1~\(\to\) #2?}
\newcommand{\msQuery}[1]{{\itshape\color{djsMuted}Query: #1}}
\newcommand{\msFlag}[2]{\textbf{#1:}\enspace #2}


\djsSchemeHeader{Complex Roots and Geometry Problems}
\djsSitePath{Core Pure}{Complex Numbers}

\djsLastUpdated{24 October 2026}
\begin{document}
\djsSetQuestionCount{13}
\djsAddMarks{1}{8}
\djsAddMarks{2}{11}
\djsAddMarks{3}{9}
\djsAddMarks{4}{10}
\djsAddMarks{5}{5}
\djsAddMarks{6}{10}
\djsAddMarks{7}{7}
\djsAddMarks{8}{12}
\djsAddMarks{9}{9}
\djsAddMarks{10}{13}
\djsAddMarks{11}{14}
\djsAddMarks{12}{8}
\djsAddMarks{13}{11}
\vspace*{6pt}
\djsContentsTitle
\vspace{30pt}
\msContentsHeading
\begin{mscontents}{}
Question & Marks by part & Marks & Page \\
\msContentsQuestion{1} & (a)(i)\,2\msSep (a)(ii)\,2\msSep (b)\,4 & 8 & \msContentsPage{1} \\
\msContentsQuestion{2} & (a)\,5\msSep (b)\,2\msSep (c)\,4 & 11 & \msContentsPage{2} \\
\msContentsQuestion{3} & (a)\,4\msSep (b)\,5 & 9 & \msContentsPage{3} \\
\msContentsQuestion{4} & (a)\,2\msSep (b)\,2\msSep (c)\,2\msSep (d)\,2\msSep (e)\,2 & 10 & \msContentsPage{4} \\
\msContentsQuestion{5} & & 5 & \msContentsPage{5} \\
\msContentsQuestion{6} & (a)\,6\msSep (b)\,4 & 10 & \msContentsPage{6} \\
\msContentsQuestion{7} & (a)\,3\msSep (b)\,4 & 7 & \msContentsPage{7} \\
\msContentsQuestion{8} & (a)\,2\msSep (b)\,1\msSep (c)\,2\msSep (d)\,3\msSep (e)\,4 & 12 & \msContentsPage{8} \\
\msContentsQuestion{9} & (a)\,2\msSep (b)\,2\msSep (c)(i)\,1\msSep (c)(ii)\,4 & 9 & \msContentsPage{9} \\
\msContentsQuestion{10} & (i)\,5\msSep (ii)\,2\msSep (iii)\,3\msSep (iv)\,3 & 13 & \msContentsPage{10} \\
\msContentsQuestion{11} & (a)(i)\,1\msSep (a)(ii)\,3\msSep (a)(iii)\,3\msSep (b)(i)\,2\msSep (b)(ii)\,2\msSep (b)(iii)\,3 & 14 & \msContentsPage{11} \\
\msContentsQuestion{12} & (a)\,2\msSep (b)\,4\msSep (c)\,2 & 8 & \msContentsPage{12} \\
\msContentsQuestion{13} & (i)\,2\msSep (ii)\,6\msSep (iii)\,1\msSep (iv)\,2 & 11 & \msContentsPage{13} \\
\SetCell[c=2]{r} Total & & 127 & \\
\end{mscontents}
\newpage
% The conventions pages, shared by every pack's mark scheme: one source of
% truth, \input after the front page, with no strand or pack logic. The keys
% on the first side, the rules on the second. The scheme layer defines
% \msHead, msList and mstblr.
\vspace*{6pt}%
{\noindent\cfttoctitlefont Reading the mark scheme\par}%
\vspace{10pt}%
These rules apply throughout, unless a note says otherwise.

\msHead{Marks}
Every mark is worth one, so \textbf{M1 A1} on a row is two marks. Only marks
coded ft follow through.\par
\vspace{10pt}%
\begin{mstblr}{
  colspec = {Q[wd=15mm, l] X[l]},
  width = \linewidth,
  hline{2} = {2.2pt, scarlet},
  column{1} = {font=\bfseries},
  row{1} = {font=\small\bfseries, rowsep=4pt},
}
Code & Meaning \\
M1 & Method: a correct method, applied to this question. A slip doesn't cost it if the method is clear. A formula must be used, not just quoted. \\
dM1 & Dependent method: as M1, but only if the method mark before it was earned, unless its note names another. \\
A1 & Accuracy: a correct answer or step. It needs the method mark it follows. \\
A1ft & As A1, following through from an earlier answer that the note names. \\
A1~AG & The last mark of a part whose answer is printed in the question: complete working, with no errors. \\
B1 & A correct value, fact or statement. It needs no method mark. \\
B1ft & As B1, following through from an earlier answer that the note names. \\
E1 & A written explanation, interpretation or conclusion. Concluding a proof by induction needs every earlier mark. \\
\end{mstblr}

\msHead{Abbreviations}
ft and AG qualify a code in the Marks column. The others appear in the
notes.\par
\vspace{10pt}%
\begin{mstblr}{
  colspec = {Q[wd=17mm, l] X[l] Q[wd=15mm, l] X[l]},
  width = \linewidth,
  column{1} = {font=\bfseries}, column{3} = {font=\bfseries},
}
ft & follow through & AG & answer given in the question \\
cao & correct answer only, with no follow-through & cso & correct solution only: no errors in the part \\
oe & or an equivalent form & awrt & any answer that rounds to this value \\
soi & seen, or implied by later working & isw & ignore working after a correct answer \\
SC & special case: credit outside the main method & their & the student's own earlier value \\
condone & a slip that doesn't cost the mark & [\;\,] & working the mark doesn't need \\
s.f. & significant figures & d.p. & decimal places \\
\end{mstblr}

\newpage
\msHead{Reading the scheme}
\begin{msList}
\item Each row is a line of working, with the marks it earns beside it.
  Where a row carries two marks, each note starts with its code.
\item \textbf{[\textit{n}]} ends each part with its total.
\item A shaded \textit{Alternative} is a second method for the same total.
  Mark an answer by the method it follows.
\item Any other valid method earns the equivalent marks, unless the question
  names a method or says \textit{hence}.
\item A sketch earns one B1 for each feature named below its figure.
\end{msList}

\msHead{Accuracy and answers}
\begin{msList}
\item Unless a question says otherwise, give answers to 3 significant
  figures, and angles in degrees to 1 decimal place. Angles in radians take
  3 significant figures.
\item A final answer with more figures is accepted if it rounds to the
  scheme's value; one with fewer is not. An answer spoiled by rounding too
  early loses its accuracy mark.
\item Where a question asks for an exact answer, a decimal loses the final
  accuracy mark. Otherwise either is accepted.
\item Wrong or missing units cost nothing unless the question asks for them.
\item A correct answer with no working earns full marks, unless the question
  needs the working: show that, prove, hence, a method it names, or an exact
  answer. A wrong answer with no working earns nothing. Calculator answers
  are treated the same way.
\item A correct answer, once reached, isn't spoiled by later working (isw).
\end{msList}

\msHead{Value of $\boldsymbol{g}$}
\begin{msList}
\item Take $g = 9.8\,\mathrm{m\,s^{-2}}$ unless a question states a value,
  which is then the only one accepted. Answers that use $g = 9.8$ may be
  given to 2 or 3 significant figures.
\item Where a question doesn't state $g$, an answer worked with $g = 10$
  throughout earns the same marks, keeping any value printed in the
  question. The worked solutions list the final answers with $g = 10$.
\end{msList}


\msQuestion{1}
\begin{msscheme}{
  hline{3-5,7-11,13-16} = {2-Z}{0.4pt, msGrid},
  hline{6,12,17} = {0.4pt, msGrid},
  cell{8-10}{2-4} = {bg=msTint},
  row{17} = {font=\bfseries},
}
Question & Scheme & Marks & Notes \\
1(a)(i) & \SetCell[c=3]{c} \msFigure[2]{\begin{tikzpicture}[scale=1]
  \coordinate (O) at (0,0);
  \coordinate (A) at (1.6,1.6);
  \coordinate (B) at (-1.6,1.6);
  \coordinate (C) at (-1.6,-1.6);
  \coordinate (D) at (1.6,-1.6);
  \draw[thick, -{Stealth[length=3mm]}] (-3.2,0) -- (3.6,0) node[right] {$\mathrm{Re}$};
  \draw[thick, -{Stealth[length=3mm]}] (0,-3.2) -- (0,3.6) node[above] {$\mathrm{Im}$};
  \fill (O) circle (1.2pt);
  \fill (A) circle (1.2pt);
  \fill (B) circle (1.2pt);
  \fill (C) circle (1.2pt);
  \fill (D) circle (1.2pt);
  \node[below left] at (O) {$O$};
  \node[above right] at (A) {$A$};
  \node[above left] at (B) {$B$};
  \node[below left] at (C) {$C$};
  \node[below right] at (D) {$D$};
\end{tikzpicture}} & & \\
& \msFeature{Labelled point} & B1 & {\(B\) and \(D\) at the quarter-turn rotations of \(A\) about \(O\), \(B\) anticlockwise from \(A\). Condone small inaccuracy in position.} \\
& \msFeature{Symmetry} & B1 & {\(C\) opposite \(A\) through \(O\), labelled \(C\).} \\
\nopagebreak & & [2] & \\
1(a)(ii) & {\(z_1 + z_1e^{\frac12\pi i} + z_1e^{\pi i} + z_1e^{\frac32\pi i}\)\\[4pt] \(\phantom{z_1} = z_1 + iz_1 - z_1 - iz_1\)} & M1 & {Evaluates the three exponentials, directly or by \(\cos\theta + i\sin\theta\), at least two correctly, and sums the resulting multiples of \(z_1\).} \\
& {\(\phantom{z_1} = z_1(1 + i - 1 - i) = 0\)} & A1~AG & {All three values correct and the terms shown to cancel to 0 without error.} \\
& \SetCell[c=3]{l} \msAlternative{Geometric series} & & \\
\nopagebreak & {\(z_1\left(1 + e^{\frac12\pi i} + e^{\pi i} + e^{\frac32\pi i}\right)\)\\[4pt] \(\phantom{z_1} = z_1\dfrac{1 - \left(e^{\frac12\pi i}\right)^4}{1 - e^{\frac12\pi i}}\)} & M1 & {Uses the sum of four terms of a geometric series with ratio \(e^{\frac12\pi i}\). Condone the factor \(z_1\) omitted here.} \\
& {\(\left(e^{\frac12\pi i}\right)^4 = e^{2\pi i} = 1 \Rightarrow\) sum \(= 0\)} & A1~AG & {Shows the ratio to the fourth power is 1 and concludes the sum is 0, with the denominator nonzero.} \\
\nopagebreak & & [2] & \\
1(b) & {\(-16 = 16e^{\pi i}\)\\[4pt] \(\phantom{-16} = 16e^{(\pi + 2k\pi)i}\)} & B1 & {Accept modulus–argument form, and the four arguments listed in place of \(\pi + 2k\pi\).} \\
& {\(z = 2e^{\frac{\pi + 2k\pi}{4}i}\), \(k = 0, 1, 2, 3\)} & M1 & {Takes the fourth root of their modulus and divides their arguments by 4, giving at least two distinct roots in any form.} \\
& {\(z = \sqrt2 + \sqrt2\,i,\ -\sqrt2 + \sqrt2\,i\)} & A1 & {Any two correct roots, in any form.} \\
& {\(z = \pm\sqrt2 \pm \sqrt2\,i\) (all four)} & A1 & {All four in the form \(a + bi\) with exact \(a\) and \(b\), and no others. Roots left in exponential form score A0.} \\
\nopagebreak & & [4] & \\
\nopagebreak \SetCell[c=2]{r} Total for Question 1 & & 8 & \\
\end{msscheme}

\msQuestion{2}
\begin{msscheme}{
  hline{3-7,9-11,13-16} = {2-Z}{0.4pt, msGrid},
  hline{8,12,17} = {0.4pt, msGrid},
  row{17} = {font=\bfseries},
}
Question & Scheme & Marks & Notes \\
2(a) & {\(|w| = \sqrt{16^2 + (16\sqrt{3})^2}\)\\[4pt] \(\arg w = \arctan\frac{16\sqrt{3}}{16}\)} & M1 & {Attempts both the modulus and the argument of \(16 + 16\sqrt{3}\,i\). Condone arithmetic slips.} \\
& {\(16 + 16\sqrt{3}\,i = 32e^{\frac{\pi}{3}i}\)} & A1 & {Or modulus 32 and argument \(\frac{\pi}{3}\) (awrt 1.05) stated.} \\
& {\(z = \text{"}32\text{"}^{\frac15}e^{\frac{\text{"}\pi/3\text{"} + 2k\pi}{5}i}\)} & M1 & {Takes the fifth root of their modulus and uses arguments (their argument \(+ 2k\pi\))\(/5\), for general \(k\) or at least two values, or rotates one root by \(\frac{2\pi}{5}\).} \\
& {\(2e^{\frac{\pi}{15}i},\ 2e^{\frac{7\pi}{15}i}, \ldots\)} & A1 & {Any two correct roots, in any form.} \\
& {\(2e^{\frac{\pi}{15}i},\ 2e^{\frac{7\pi}{15}i},\ 2e^{\frac{13\pi}{15}i},\)\\[4pt] \(2e^{\frac{19\pi}{15}i},\ 2e^{\frac{5\pi}{3}i}\)} & A1 & {All five as \(2e^{i\theta}\) with \(0 \leq \theta < 2\pi\), exact or awrt 0.209, 1.47, 2.72, 3.98, 5.24, and no others. A0 for any argument outside the range.} \\
\nopagebreak & & [5] & \\
2(b) & \SetCell[c=3]{c} \msFigure{\begin{tikzpicture}
\begin{axis}[
axis lines=middle,
xlabel={$\mathrm{Re}$},
ylabel={$\mathrm{Im}$},
xmin=-3.2, xmax=3.2,
ymin=-3.2, ymax=3.2,
xtick={-2,2},
ytick={-2,2},
width=11cm,
height=11cm,
axis equal image,
clip=false
]
\addplot[black, thick] coordinates {
(1.956,0.416)
(0.209,1.989)
(-1.827,0.813)
(-1.338,-1.486)
(1.000,-1.732)
(1.956,0.416)
};
\addplot[only marks, mark=*] coordinates {
(1.956,0.416)
(0.209,1.989)
(-1.827,0.813)
(-1.338,-1.486)
(1.000,-1.732)
};
\node[anchor=west] at (axis cs:1.956,0.416) {$A$};
\node[anchor=south] at (axis cs:0.209,1.989) {$B$};
\node[anchor=east] at (axis cs:-1.827,0.813) {$C$};
\node[anchor=north east] at (axis cs:-1.338,-1.486) {$D$};
\node[anchor=west] at (axis cs:1.000,-1.732) {$E$};
\node[anchor=north east] at (axis cs:0,0) {$O$};
\end{axis}
\end{tikzpicture}} & & \\
& \msFeature{Shape} & B1 & {Pentagon centred at \(O\) with vertices equally spaced on a circle of radius 2, the radius shown by a scale or label.} \\
& \msFeature{Labelled point} & B1 & {Each vertex labelled with its root, or a letter identified with it, in its correct position, one near the positive real axis in the first quadrant.} \\
\nopagebreak & & [2] & \\
2(c) & {Five congruent triangles with sides \(2, 2\)\\[4pt] and angle \(\frac{2\pi}{5}\) at \(O\)} & B1 & {Angle at \(O\) may be implied by its use in an area.} \\
& {\(\frac12 \times 2 \times 2 \times \sin\frac{2\pi}{5}\)} & M1 & {Uses \(\frac12 ab\sin C\) with two sides equal to their modulus and their angle at \(O\) as the included angle.} \\
& {\(\phantom{\tfrac12} = 2\sin\frac{2\pi}{5}\)} & A1 & \\
& {\(5 \times 2\sin\frac{2\pi}{5} = 10\sin\frac{2\pi}{5} \Rightarrow k = 10\)} & A1 & {cso. Area \(10\sin\frac{2\pi}{5}\) or \(k = 10\) stated; a decimal area alone scores A0.} \\
\nopagebreak & & [4] & \\
\nopagebreak \SetCell[c=2]{r} Total for Question 2 & & 11 & \\
\end{msscheme}

\msQuestion{3}
\begin{msscheme}{
  hline{3-6,8-11} = {2-Z}{0.4pt, msGrid},
  hline{7,12} = {0.4pt, msGrid},
  row{12} = {font=\bfseries},
}
Question & Scheme & Marks & Notes \\
3(a) & {\((h-1)^2 + 4^2 = (h-7)^2 + 2^2\)\\[4pt] \(\Rightarrow h = \ldots\)} & M1 & {Equates the distances, or squared distances, from \((h, 0)\) to \(A\) and \(B\), or intersects the perpendicular bisector of \(AB\) with the real axis. Condone sign slips.} \\
& {\(h = 3\)} & A1 & \\
& {\(r^2 = (\text{"}3\text{"} - 1)^2 + 4^2\ [= 20]\)} & M1 & {Uses Pythagoras for the distance from their real centre to \(A\) or \(B\).} \\
& {\(|z - 3| = 2\sqrt{5}\)} & A1 & {Accept \(\sqrt{20}\) or awrt 4.47. Only the form \(|z-a|=b\) scores; \((x-3)^2 + y^2 = 20\) alone scores A0.} \\
\nopagebreak & & [4] & \\
3(b) & {\((x-3)^2 + y^2 = 20\)\\[4pt] \((x-1)^2 + (y-1)^2 = 25\)} & M1 A1 & {\msNote{M1}{Writes at least one circle as an equation in \(x\) and \(y\), using \(z = x + iy\) and their \(C\) of the form \(|z-a|=b\). Condone sign slips.}\\ \msNote{A1}{Both correct, in any form.}} \\
& {\(4x - 2y - 12 = 0 \Rightarrow y = 2x - 6\)} & M1 & {Subtracts their equations to eliminate \(x^2\) and \(y^2\), reaching a linear equation in \(x\) and \(y\). Condone sign slips.} \\
& {\((x-3)^2 + (\text{"}2x-6\text{"})^2 = 20\)\\[4pt] \(\Rightarrow x = \ldots\)} & dM1 & {Substitutes their linear equation into either circle and solves the resulting quadratic in one variable by any method.} \\
& {\(z_1 = 1 - 4i,\ z_2 = 5 + 4i\)} & A1 & {Both, in either order, as complex numbers; coordinates alone score A0.} \\
\nopagebreak & & [5] & \\
\nopagebreak \SetCell[c=2]{r} Total for Question 3 & & 9 & \\
\end{msscheme}

\msQuestion{4}
\begin{msscheme}{
  hline{3-4,6-8,10-11,13-14,16-17} = {2-Z}{0.4pt, msGrid},
  hline{5,9,12,15,18} = {0.4pt, msGrid},
  row{18} = {font=\bfseries},
}
Question & Scheme & Marks & Notes \\
4(a) & {\(z_2 = \frac{(7 - i)(1 + i)}{(1 - i)(1 + i)} = \ldots\)} & M1 & {Multiplies numerator and denominator by \(1 + i\) and expands, or equates real and imaginary parts of \((a + bi)(1 - i) = 7 - i\). Condone one slip.} \\
& {\(z_2 = 4 + 3i\)} & A1 & \\
\nopagebreak & & [2] & \\
4(b) & \SetCell[c=3]{c} \msFigure{\begin{tikzpicture}
\begin{axis}[
  axis lines=middle,
  xmin=-0.5, xmax=9,
  ymin=-2, ymax=5,
  width=10cm, height=7.5cm,
  xlabel={$\mathrm{Re}$},
  ylabel={$\mathrm{Im}$},
  xtick={4,7},
  ytick={-1,3},
  enlargelimits=false
]
\addplot[only marks, mark=*, mark size=2.2pt] coordinates {(7,-1) (4,3)};
\node at (axis cs:7.25,-1.2) {$P$};
\node at (axis cs:4.25,3.2) {$Q$};
\end{axis}
\end{tikzpicture}} & & \\
& \msFeature{Labelled point} & B1 & {\(P\) plotted and labelled at \((7, -1)\), in the fourth quadrant, with its position shown by scale or coordinates.} \\
& \msFeature{Labelled point} & B1ft & {ft their \(z_2\) from (a). \(Q\) plotted and labelled at their point, positioned correctly relative to \(P\).} \\
\nopagebreak & & [2] & \\
4(c) & {\(\overrightarrow{QO}\cdot\overrightarrow{QP} = (-4)(3) + (-3)(-4)\)} & M1 & {Applies a perpendicularity test at \(Q\) to \(OQ\) and \(QP\), such as a scalar product, a product of gradients or Pythagoras' converse. Condone sign slips.} \\
& {\(\phantom{\overrightarrow{QO}\cdot\overrightarrow{QP}} = 0 \Rightarrow \angle OQP = \frac{\pi}{2}\)} & A1~AG & {Correct values giving 0, a gradient product of \(-1\), or \(OQ^2 + QP^2 = OP^2\), with a conclusion that the angle is a right angle.} \\
\nopagebreak & & [2] & \\
4(d) & {Uses \(OP\) as a diameter: \(C = \tfrac12(7 - i) = \ldots\)} & M1 & {Takes \(C\) as the midpoint of \(OP\), or solves for the point equidistant from \(O\), \(P\) and \(Q\).} \\
& {\(C = \tfrac72 - \tfrac12i\)} & A1 & {oe as a complex number, such as \(3.5 - 0.5i\). Coordinates \(\left(\frac72, -\frac12\right)\) alone do not score.} \\
\nopagebreak & & [2] & \\
4(e) & {\(r = \tfrac12|OP| = \tfrac12\sqrt{7^2 + (-1)^2}\)} & M1 & {Finds half of \(|OP|\), or the distance from their \(C\) to \(O\), \(P\) or \(Q\), by Pythagoras. Condone sign slips.} \\
& {\(r = \frac{5\sqrt2}{2}\)} & A1 & {Any exact form, such as \(\frac{\sqrt{50}}{2}\) or \(\sqrt{12.5}\).} \\
\nopagebreak & & [2] & \\
\nopagebreak \SetCell[c=2]{r} Total for Question 4 & & 10 & \\
\end{msscheme}

\msQuestion{5}
\begin{msscheme}{
  hline{3-6} = {2-Z}{0.4pt, msGrid},
  hline{7} = {0.4pt, msGrid},
  row{7} = {font=\bfseries},
}
Question & Scheme & Marks & Notes \\
5 & {\(A\), \(B\), \(C\): \(z^3 - 1 = 0\)} & B1 & \\
& {\(P\), \(Q\), \(R\): \((z - 2)^3 = 1\)} & M1 & {Replaces \(z\) by \(z - 2\) in their cubic, or forms a cubic from the three roots by factors or root sums and products. Need not be expanded.} \\
& {\(z^3 - 6z^2 + 12z - 9 = 0\)} & A1 & {Correct cubic, oe such as \((z - 3)(z^2 - 3z + 3) = 0\).} \\
& {\((z^3 - 1)(z^3 - 6z^2 + 12z - 9) = 0\)\\[4pt] \(\phantom{(z^3 - 1)} \Rightarrow z^6 - 6z^5 + \ldots = 0\)} & M1 & {Multiplies their two cubics and expands. Condone slips in collecting terms.} \\
& {\(z^6 - 6z^5 + 12z^4 - 10z^3 + 6z^2 - 12z + 9 = 0\)} & A1 & {Any nonzero integer multiple accepted.} \\
\nopagebreak \SetCell[c=2]{r} Total for Question 5 & & 5 & \\
\end{msscheme}

\msQuestion{6}
\begin{msscheme}{
  hline{3-8,10-13} = {2-Z}{0.4pt, msGrid},
  hline{9,14} = {0.4pt, msGrid},
  row{14} = {font=\bfseries},
}
Question & Scheme & Marks & Notes \\
6(a) & {\(z_A - z_G = 4 + 4i\)} & B1 & \\
& {\(z_B - z_G = (4+4i)\left(-\tfrac12 + \tfrac{\sqrt3}{2}i\right)\)\\[4pt] \(\phantom{z_B - z_G} = \ldots\)} & M1 & {Multiplies their \(z_A - z_G\) by \(e^{\pm 2\pi i/3}\) in any form and expands. Condone a sign slip in the cosine or sine value.} \\
& {\(z_B - z_G = (-2-2\sqrt3) + (2\sqrt3-2)i\)} & A1 & \\
& {\(z_B = (-1-2\sqrt3) + (2\sqrt3-4)i\)} & A1 & {Exact, in the form \(a + bi\). \(B\) and \(C\) may be interchanged.} \\
& {\(z_C = (1-2i) + (4+4i)\left(-\tfrac12 - \tfrac{\sqrt3}{2}i\right)\)\\[4pt] \(\phantom{z_C} = \ldots\)} & M1 & {Rotates their \(z_A - z_G\) through the other angle and adds \(z_G\), or uses \(z_C = 3z_G - z_A - z_B\) with their \(z_B\).} \\
& {\(z_C = (-1+2\sqrt3) - (4+2\sqrt3)i\)} & A1 & {Exact, in the form \(a + bi\).} \\
\nopagebreak & & [6] & \\
6(b) & {\(z_D = \tfrac13(z_G + z_B + z_C)\), \(z_B + z_C = 3z_G - z_A\)\\[4pt] \(\Rightarrow z_D = \tfrac13(4z_G - z_A)\)} & M1 & {Writes \(z_D\) as the mean of \(z_G\), \(z_B\), \(z_C\), then simplifies using \(z_A + z_B + z_C = 3z_G\) or evaluates with their (a) values.} \\
& {\(z_D - z_G = -\tfrac13(z_A - z_G)\), similarly for \(E\), \(F\)\\[4pt] \(\Rightarrow\) area factor \(\left(\tfrac13\right)^2 = \tfrac19\)} & A1 & {Lengths scaled by \(\frac13\) about \(G\), sign not needed, with \(E\) and \(F\) by symmetry, and area factor \(\frac19\).} \\
& {\(GA = |4 + 4i| = 4\sqrt2\)\\[4pt] Area \(ABC = 3 \times \tfrac12(4\sqrt2)^2\sin\tfrac{2\pi}{3}\)\\[4pt] \(\phantom{\text{Area } ABC} = \ldots\ [= 24\sqrt3]\)} & M1 & {Uses three triangles \(\frac12 r^2 \sin\frac{2\pi}{3}\) with their \(r = |z_A - z_G|\), or \(\frac{\sqrt3}{4}s^2\) with their side \(s = |z_A - z_B|\).} \\
& {Area \(DEF = \tfrac19 \times 24\sqrt3 = \tfrac{8\sqrt3}{3}\)} & A1 & {Exact, oe such as \(\frac{8}{\sqrt3}\).} \\
\nopagebreak & & [4] & \\
\nopagebreak \SetCell[c=2]{r} Total for Question 6 & & 10 & \\
\end{msscheme}

\msQuestion{7}
\begin{msscheme}{
  hline{3-5,7-10} = {2-Z}{0.4pt, msGrid},
  hline{6,11} = {0.4pt, msGrid},
  row{11} = {font=\bfseries},
}
Question & Scheme & Marks & Notes \\
7(a) & {\(|w-(1+i)|^2 = \bigl(w-(1+i)\bigr)\bigl(w^*-(1-i)\bigr)\)\\[4pt] \(\phantom{|w-(1+i)|^2} = \ldots\)} & M1 & {Multiplies \(w-(1+i)\) by its attempted conjugate and expands, or an equivalent Cartesian method. Condone sign slips.} \\
& {\(= ww^* - (1-i)w - (1+i)w^* + (1+i)(1-i)\)} & A1 & {Correct four-term expansion, oe.} \\
& {\(ww^* = |w|^2 = 1\), \((1+i)(1-i) = 2\)\\[4pt] \(\Rightarrow |w-(1+i)|^2 = 3-(1-i)w-(1+i)w^*\)} & A1~AG & {Uses \(ww^* = 1\) from \(|w| = 1\) and \((1+i)(1-i) = 2\), reaching the printed result without error.} \\
\nopagebreak & & [3] & \\
7(b) & {\(A_{2r-1}\colon w_r = e^{\frac{2(r-1)\pi}{5}i}\), the fifth roots of unity\\[4pt] \(P\colon 1+i\)} & B1 & {Needs both the odd-numbered vertices as the fifth roots of unity, or their five arguments, and \(P\) as \(1+i\).} \\
& {\(\sum_{r=1}^{5}(PA_{2r-1})^2 = \sum_{r=1}^{5}\bigl(3-(1-i)w_r-(1+i)w_r^*\bigr)\)\\[4pt] \(\phantom{\sum} = 15-(1-i)\sum w_r-(1+i)\sum w_r^*\)} & M1 & {Sums the result of (a) over the five odd-numbered vertices, separating the sums of \(w_r\) and \(w_r^*\). Condone the constant left as \(\sum 3\).} \\
& {\(\sum w_r = 0\), \(\sum w_r^* = 0\)} & B1 & {Both sums stated or shown to be zero.} \\
& {\(\sum_{r=1}^{5}(PA_{2r-1})^2 = 15\)} & A1 & {cso.} \\
\nopagebreak & & [4] & \\
\nopagebreak \SetCell[c=2]{r} Total for Question 7 & & 7 & \\
\end{msscheme}

\msQuestion{8}
\begin{msscheme}{
  hline{3-4,6,8-9,11-13,15-18} = {2-Z}{0.4pt, msGrid},
  hline{5,7,10,14,19} = {0.4pt, msGrid},
  row{19} = {font=\bfseries},
}
Question & Scheme & Marks & Notes \\
8(a) & {\(\alpha^5 = \cos\left(5 \times \frac{3\pi}{5}\right) + i\sin\left(5 \times \frac{3\pi}{5}\right)\)\\[4pt] \(\phantom{\alpha^5} = \cos 3\pi + i\sin 3\pi\)} & M1 & {Applies de Moivre's theorem to \(\alpha^5\), reaching argument \(3\pi\) or an equivalent.} \\
& {\(\alpha^5 = -1\), so \(\alpha\) is a root} & A1~AG & {Evaluates \(\cos 3\pi = -1\) and \(\sin 3\pi = 0\) and concludes that \(\alpha\) is a root.} \\
\nopagebreak & & [2] & \\
8(b) & {\(\cos\frac{\pi}{5} + i\sin\frac{\pi}{5}\), \(\ -1\),\\[4pt] \(\cos\frac{7\pi}{5} + i\sin\frac{7\pi}{5}\), \(\ \cos\frac{9\pi}{5} + i\sin\frac{9\pi}{5}\)} & B1 & {All four, in any correct form, with arguments in any range such as \(-\pi < \theta \leq \pi\).} \\
\nopagebreak & & [1] & \\
8(c) & {\(\alpha^5 + 1 = (\alpha + 1)(\alpha^4 - \alpha^3 + \alpha^2 - \alpha + 1)\)} & M1 & {Factorises \(z^5 + 1\) with factor \(z + 1\), or sums the geometric series with ratio \(-\alpha\). Expanding the printed product to \(\alpha^5 + 1\) also scores.} \\
& {\(\alpha^5 + 1 = 0\) and \(\alpha \neq -1\)\\[4pt] \(\Rightarrow \alpha^4 - \alpha^3 + \alpha^2 - \alpha + 1 = 0\)} & A1~AG & {Uses \(\alpha^5 + 1 = 0\) and states \(\alpha \neq -1\) before concluding. A0 if \(\alpha \neq -1\) is not stated.} \\
\nopagebreak & & [2] & \\
8(d) & {\(\alpha^2 - \alpha + 1 - \frac{1}{\alpha} + \frac{1}{\alpha^2} = 0\)} & M1 & {Divides the quartic by \(\alpha^2\), or multiplies an expansion of the printed LHS by \(\alpha^2\). Condone one slip.} \\
& {\(\left(\alpha + \frac{1}{\alpha}\right)^2 = \alpha^2 + 2 + \frac{1}{\alpha^2}\)\\[4pt] \(\Rightarrow \left(\alpha + \frac{1}{\alpha}\right)^2 - 2 - \left(\alpha + \frac{1}{\alpha}\right) + 1 = 0\)} & M1 & {Uses the expansion of \(\left(\alpha + \frac{1}{\alpha}\right)^2\) to replace \(\alpha^2 + \frac{1}{\alpha^2}\), or to expand the printed LHS.} \\
& {\(\left(\alpha + \frac{1}{\alpha}\right)^2 - \left(\alpha + \frac{1}{\alpha}\right) - 1 = 0\)} & A1~AG & {Printed result reached without error from the result of (c). Showing the printed LHS equals \(\alpha^{-2}\) times the quartic, then citing (c), also scores.} \\
\nopagebreak & & [3] & \\
8(e) & {\(t^2 - t - 1 = 0 \Rightarrow t = \ldots\), where \(t = \alpha + \frac{1}{\alpha}\)} & M1 & {Solves the quadratic in \(\alpha + \frac{1}{\alpha}\) by formula, completing the square or calculator.} \\
& {\(t = \frac{1 \pm \sqrt{5}}{2}\)} & A1 & \\
& {\(\frac{1}{\alpha} = \cos\frac{3\pi}{5} - i\sin\frac{3\pi}{5}\)\\[4pt] \(\Rightarrow 2\cos\frac{3\pi}{5} = \text{"}\frac{1 \pm \sqrt{5}}{2}\text{"} \Rightarrow \cos\frac{3\pi}{5} = \ldots\)} & M1 & {Shows or quotes \(\alpha + \frac{1}{\alpha} = 2\cos\frac{3\pi}{5}\) and equates it to their roots.} \\
& {\(\cos\frac{3\pi}{5} < 0 \Rightarrow \cos\frac{3\pi}{5} = \frac14 - \frac14\sqrt{5}\)} & A1 & {Positive value rejected as \(\cos\frac{3\pi}{5} < 0\). In the form \(a + b\sqrt{c}\), or \(a\), \(b\), \(c\) stated; awrt \(-0.309\) scores A0.} \\
\nopagebreak & & [4] & \\
\nopagebreak \SetCell[c=2]{r} Total for Question 8 & & 12 & \\
\end{msscheme}

\msQuestion{9}
\begin{msscheme}{
  hline{3-4,6-7,9,11-14} = {2-Z}{0.4pt, msGrid},
  hline{5,8,10,15} = {0.4pt, msGrid},
  row{15} = {font=\bfseries},
}
Question & Scheme & Marks & Notes \\
9(a) & {\(\alpha = \frac{(6+11i) + (-8-3i)}{2}\)} & M1 & {Uses the midpoint of \(z_1\) and \(z_2\), with values substituted.} \\
& {\(\alpha = \frac{-2+8i}{2} = -1+4i\)} & A1~AG & {Printed result reached from the substituted midpoint without error.} \\
\nopagebreak & & [2] & \\
9(b) & {\(z_1 - \alpha = 7+7i\)\\[4pt] \(\beta(z_1 - \alpha) = \frac{(1-i)(7+7i)}{14}\)} & M1 & {Finds \(z_1 - \alpha\) and multiplies by \(\beta\). Condone a slip in \(z_1 - \alpha\).} \\
& {\(\phantom{\beta(z_1 - \alpha)} = \frac{7(1-i)(1+i)}{14} = \frac{7 \times 2}{14} = 1\)} & A1~AG & {Product evaluated explicitly, without error, to 1; a bare calculator result scores A0.} \\
\nopagebreak & & [2] & \\
9(c)(i) & {\(w = e^{\frac{k\pi}{4}i}, \quad k = 0, 1, \ldots, 7\)} & B1 & {All eight roots, for any eight consecutive \(k\) or listed separately. \(r = 1\) may be omitted.} \\
\nopagebreak & & [1] & \\
9(c)(ii) & {\(w = \beta(z - \alpha) \Rightarrow z = \alpha + \frac{w}{\beta}\)} & M1 & {Rearranges for \(z\) in terms of \(w\), or otherwise rotates \(z_1 - \alpha\) about \(\alpha\) through multiples of \(\frac{\pi}{4}\).} \\
& {\(\frac{1}{\beta} = 7(1+i)\)\\[4pt] \(z = -1+4i+7(1+i)w\)} & A1 & {oe} \\
& {\(w = \frac{1+i}{\sqrt2} \Rightarrow z = -1+4i+7(1+i)\frac{1+i}{\sqrt2}\)\\[4pt] \(w = i \Rightarrow z = -1+4i+7(1+i)i\), \(\ldots\)} & dM1 & {Substitutes at least two of the six non-real roots into their \(z = \alpha + cw\), for a constant \(c\).} \\
& {\(-1+(4+7\sqrt2)i,\ -8+11i,\ (-1-7\sqrt2)+4i\)\\[4pt] \(-1+(4-7\sqrt2)i,\ 6-3i,\ (-1+7\sqrt2)+4i\)} & A1 & {All six, exact or awrt 3 s.f., and no others; \(z_1\) and \(z_2\) also listed is condoned.} \\
\nopagebreak & & [4] & \\
\nopagebreak \SetCell[c=2]{r} Total for Question 9 & & 9 & \\
\end{msscheme}

\msQuestion{10}
\begin{msscheme}{
  hline{3-7,9-11,13-15,17-19} = {2-Z}{0.4pt, msGrid},
  hline{8,12,16,20} = {0.4pt, msGrid},
  row{20} = {font=\bfseries},
}
Question & Scheme & Marks & Notes \\
10(i) & {\(|z| = \sqrt{128^2 + (128\sqrt3)^2}\)\\[4pt] \(\arg z = \pi - \arctan\sqrt3\)} & M1 & {Attempts the modulus and an argument adjusted for the second quadrant. Condone arithmetic slips.} \\
& {\(-128 + 128\sqrt3\,i = 256e^{\frac{2\pi}{3}i}\)} & A1 & \\
& {\(z_k = 256^{\frac18}e^{\frac{1}{8}\left(\frac{2\pi}{3} + 2k\pi\right)i}\)\\[4pt] \(\phantom{z_k} = 2e^{\left(\frac{\pi}{12} + \frac{k\pi}{4}\right)i}\)} & M1 & {Takes \(r^{\frac18}\) and \(\frac{\theta + 2k\pi}{8}\) from their \(re^{i\theta}\), for at least two values of \(k\).} \\
& {Two of \(2e^{\frac{\pi}{12}i}, 2e^{\frac{\pi}{3}i}, \ldots\)} & A1 & {Any two correct roots, arguments in any range.} \\
& {\(2e^{\frac{\pi}{12}i}, 2e^{\frac{\pi}{3}i}, 2e^{\frac{7\pi}{12}i}, 2e^{\frac{5\pi}{6}i}\)\\[4pt] \(2e^{\frac{13\pi}{12}i}, 2e^{\frac{4\pi}{3}i}, 2e^{\frac{19\pi}{12}i}, 2e^{\frac{11\pi}{6}i}\)} & A1 & {All eight with \(r = 2\) and arguments in \(0 \leq \theta < 2\pi\), and no others.} \\
\nopagebreak & & [5] & \\
10(ii) & \SetCell[c=3]{c} \msFigure{\begin{tikzpicture}
\begin{axis}[
axis lines=middle,
xlabel={$\mathrm{Re}$},
ylabel={$\mathrm{Im}$},
xmin=-3.5, xmax=3.5,
ymin=-3.5, ymax=3.5,
width=12cm, height=12cm,
axis equal image,
ticks=none,
clip=false
]
\addplot[only marks, mark=*, mark size=1.6pt] coordinates {
(1.9319,0.5176)
(1,1.7321)
(-0.5176,1.9319)
(-1.7321,1)
(-1.9319,-0.5176)
(-1,-1.7321)
(0.5176,-1.9319)
(1.7321,-1)
};
\node[anchor=west] at (axis cs:1.9319,0.5176) {$A$};
\node[anchor=west] at (axis cs:1,1.7321) {$B$};
\node[anchor=south east] at (axis cs:-0.5176,1.9319) {$C$};
\node[anchor=east] at (axis cs:-1.7321,1) {$D$};
\node[anchor=east] at (axis cs:-1.9319,-0.5176) {$E$};
\node[anchor=east] at (axis cs:-1,-1.7321) {$F$};
\node[anchor=north east] at (axis cs:0.5176,-1.9319) {$G$};
\node[anchor=west] at (axis cs:1.7321,-1) {$H$};
\end{axis}
\end{tikzpicture}} & & \\
& \msFeature{Shape} & B1 & {Eight points equally spaced on a circle centred at \(O\).} \\
& \msFeature{Labelled point} & B1 & {\(A\) to \(H\) anticlockwise, \(A\) at argument \(\frac{\pi}{12}\) in the first quadrant.} \\
\nopagebreak & & [2] & \\
10(iii) & {\(w = \tfrac12\left(2e^{\frac{\pi}{3}i} + 2e^{\frac{5\pi}{6}i}\right)\)} & M1 & {Halves the sum of their second and fourth roots, in any form; Cartesian accepted.} \\
& {\(w = e^{\frac{7\pi}{12}i}\left(e^{-\frac{\pi}{4}i} + e^{\frac{\pi}{4}i}\right)\)\\[4pt] \(|w| = \sqrt2\)} & A1 & {\(\sqrt2\) or an exact equivalent.} \\
& {\(\arg w = \frac{7\pi}{12}\)} & A1 & {Or an exact equivalent.} \\
\nopagebreak & & [3] & \\
10(iv) & {\(w^n = (\sqrt2)^ne^{\frac{7n\pi}{12}i}\)\\[4pt] \(\frac{7n\pi}{12} = m\pi\), \(m \in \mathbb{Z}\)} & M1 & {Requires \(n\) times their \(\arg w\) to be an integer multiple of \(\pi\). Implied by \(n = 12\).} \\
& {\(n = 12\)} & A1 & \\
& {\(a = (\sqrt2)^{12}e^{7\pi i} = -64\)} & A1 & \\
\nopagebreak & & [3] & \\
\nopagebreak \SetCell[c=2]{r} Total for Question 10 & & 13 & \\
\end{msscheme}

\msQuestion{11}
\begin{msscheme}{
  hline{3,5-7,9-11,13-14,16-17,19-21} = {2-Z}{0.4pt, msGrid},
  hline{4,8,12,15,18,22} = {0.4pt, msGrid},
  row{22} = {font=\bfseries},
}
Question & Scheme & Marks & Notes \\
11(a)(i) & {\(\beta = -2 - 2i\)} & B1 & \\
\nopagebreak & & [1] & \\
11(a)(ii) & {\(\alpha\beta = (-2+2i)(-2-2i) = 8\)} & B1 & {Or the quadratic factor \(z^2 + 4z + 8\).} \\
& {\(\text{"}8\text{"}\gamma = -24 \Rightarrow \gamma = \ldots\)} & M1 & {Sets \(\alpha\beta\gamma = -24\) with their \(\alpha\beta\), or divides the cubic by their quadratic factor or compares constant terms, to find \(\gamma\).} \\
& {\(\gamma = -3\)} & A1 & \\
\nopagebreak & & [3] & \\
11(a)(iii) & {\(-p = -4 + \text{"}{-3}\text{"} \Rightarrow p = \ldots\)\\[4pt] \(q = 8 + \text{"}{-3}\text{"}(-4) \Rightarrow q = \ldots\)} & M1 & {Uses \(\sum\alpha = -p\) or \(\sum\alpha\beta = q\) with their roots, or expands their factors \((z^2+4z+8)(z - \gamma)\) and compares coefficients.} \\
& {\(p = 7\)} & A1 & \\
& {\(q = 20\)} & A1 & \\
\nopagebreak & & [3] & \\
11(b)(i) & {\(r = \sqrt{(-2)^2 + 2^2}\), \(\theta = \pi - \arctan 1\)} & M1 & {Attempts \(|\alpha|\) and an argument in the second quadrant. Not for \(\pm\frac{\pi}{4}\).} \\
& {\(\alpha = 2\sqrt{2}\,e^{\frac{3\pi}{4}i}\)} & A1 & {\(r = 2\sqrt{2}\), \(\sqrt{8}\) or awrt 2.83 and \(\theta = \frac{3\pi}{4}\) or awrt 2.36, in the form \(re^{i\theta}\); not for \(r\) and \(\theta\) stated alone.} \\
\nopagebreak & & [2] & \\
11(b)(ii) & {\((-2+2i)^n = \left(2\sqrt{2}\,e^{\frac{3\pi}{4}i}\right)^n\)\\[4pt] \(\phantom{(-2+2i)^n} = (2\sqrt{2})^ne^{\frac{3n\pi}{4}i}\)} & M1 & {Applies \((re^{i\theta})^n = r^ne^{in\theta}\), or de Moivre, to their \(\alpha\) in modulus-argument form from (b)(i).} \\
& {\(\phantom{(-2+2i)^n} = (2\sqrt{2})^n\left(\cos\frac{3n\pi}{4} + i\sin\frac{3n\pi}{4}\right)\)} & A1~AG & {Correct \(\alpha^n\) from \(2\sqrt{2}\,e^{\frac{3\pi}{4}i}\) leading without error to the printed result; combined steps accepted.} \\
\nopagebreak & & [2] & \\
11(b)(iii) & {\(\beta^n = \left(2\sqrt{2}\,e^{-\frac{3\pi}{4}i}\right)^n\)\\[4pt] \(\phantom{\beta^n} = (2\sqrt{2})^n\left(\cos\frac{3n\pi}{4} - i\sin\frac{3n\pi}{4}\right)\)} & M1 & {Raises \(\beta\) in exponential or modulus-argument form, argument \(-\frac{3\pi}{4}\) oe, to the power \(n\), or takes \(\beta^n\) as the conjugate of \(\alpha^n\).} \\
& {\(\alpha^n + \beta^n = 2(2\sqrt{2})^n\cos\frac{3n\pi}{4}\)} & A1 & \\
& {\(\gamma^n = (-3)^n\)\\[4pt] \(\alpha^n + \beta^n + \gamma^n\)\\[4pt] \(\phantom{\alpha^n} = 2(2\sqrt{2})^n\cos\frac{3n\pi}{4} + (-3)^n\)} & A1~AG & {Adds \((-3)^n\) from \(\gamma = -3\) to a correctly obtained \(\alpha^n + \beta^n\), reaching the printed result without error.} \\
\nopagebreak & & [3] & \\
\nopagebreak \SetCell[c=2]{r} Total for Question 11 & & 14 & \\
\end{msscheme}

\msQuestion{12}
\begin{msscheme}{
  hline{3-4,6-9,11-12} = {2-Z}{0.4pt, msGrid},
  hline{5,10,13} = {0.4pt, msGrid},
  row{13} = {font=\bfseries},
}
Question & Scheme & Marks & Notes \\
12(a) & {\(-729 = 3^6e^{(\pi + 2k\pi)i}\)\\[4pt] \(\Rightarrow z = 3e^{\left(\frac{\pi}{6} + \frac{k\pi}{3}\right)i}\)} & B1 & {Modulus 3 and arguments \(\frac{\pi}{6} + \frac{k\pi}{3}\), soi by at least two correct roots.} \\
& {\(A: 3e^{\frac{\pi}{6}i}\), \(B: 3e^{\frac{\pi}{2}i}\), \(C: 3e^{\frac{5\pi}{6}i}\)\\[4pt] \(D: 3e^{\frac{7\pi}{6}i}\), \(E: 3e^{\frac{3\pi}{2}i}\), \(F: 3e^{\frac{11\pi}{6}i}\)} & B1 & {All six in exponential form; arguments may differ by multiples of \(2\pi\), such as \(-\frac{\pi}{6}\) for \(F\).} \\
\nopagebreak & & [2] & \\
12(b) & {\(z_G = \frac{z_A + z_B}{3}\), \(w = \frac{z_G}{z_A} = \frac{z_A + z_B}{3z_A}\)} & M1 & {Forms a centroid as the mean of its three vertices, with \(O\) as 0, and divides it by the corresponding root.} \\
& {\(z_B = z_Ae^{\frac{\pi}{3}i} \Rightarrow w = \frac{1 + e^{\frac{\pi}{3}i}}{3}\)} & A1 & {Any correct expression for \(w\), oe such as \(\frac12 + \frac{\sqrt3}{6}i\).} \\
& {\(1 + e^{\frac{\pi}{3}i} = \frac32 + \frac{\sqrt3}{2}i\)\\[4pt] \(\Rightarrow r = \ldots\ [= \sqrt3]\), \(\theta = \ldots\ [= \frac{\pi}{6}]\)} & M1 & {Finds the modulus and argument of their \(w\) or its numerator, by any method such as Cartesian form or factorising out \(e^{\frac{\pi}{6}i}\). Condone slips.} \\
& {\(w = \frac{1}{\sqrt3}e^{\frac{\pi}{6}i}\)} & A1 & {Exponential form required; oe such as \(\frac{\sqrt3}{3}e^{\frac{\pi}{6}i}\), or awrt 0.577 and 0.524.} \\
\nopagebreak & & [4] & \\
12(c) & {\(p = \text{"}w\text{"}^6 \times (-729)\)\\[4pt] or \(p = z_G^{\,6} = \left(\sqrt3\,e^{\frac{\pi}{3}i}\right)^6\)} & M1 & {Multiplies the sixth power of their \(w\) by \(-729\), or finds the sixth power of one of their points \(G\) to \(L\).} \\
& {\(p = 27\)} & A1 & \\
\nopagebreak & & [2] & \\
\nopagebreak \SetCell[c=2]{r} Total for Question 12 & & 8 & \\
\end{msscheme}

\msQuestion{13}
\begin{msscheme}{
  hline{3,5-10,12,14-15} = {2-Z}{0.4pt, msGrid},
  hline{4,11,13,16} = {0.4pt, msGrid},
  row{16} = {font=\bfseries},
}
Question & Scheme & Marks & Notes \\
13(i) & {\(w = e^{\frac{(2k+1)\pi}{11}i}\)\\[4pt] \(k = 0, 1, \ldots, 10\)} & B1 B1 & {\msNote{B1}{Any correct general form, such as modulus-argument; \(k\) need not be restricted.}\\ \msNote{B1}{Exactly eleven distinct roots, by a range of \(k\) such as \(-5\) to \(5\), or listed.}} \\
\nopagebreak & & [2] & \\
13(ii) & {\(\left(\frac{z+1}{z-1}\right)^{11} = -1\)\\[4pt] \(\Rightarrow \frac{z+1}{z-1} = w = e^{i\theta}\), \(\theta = \frac{(2k+1)\pi}{11}\)} & M1 & {Divides by \((z-1)^{11}\) and sets \(\frac{z+1}{z-1}\) equal to their roots from (i). Condone \(z \neq 1\) not stated.} \\
& {\(w(z-1) = z+1 \Rightarrow z = \frac{w+1}{w-1}\)} & M1 & {Rearranges to make \(z\) the subject. Condone sign slips.} \\
& {\(z = \frac{e^{i\theta/2}+e^{-i\theta/2}}{e^{i\theta/2}-e^{-i\theta/2}}\)\\[4pt] \(\phantom{z} = \frac{2\cos(\theta/2)}{2i\sin(\theta/2)}\)} & dM1 & {Multiplies by \(e^{-i\theta/2}\), or uses half-angle identities on \(1 \pm \cos\theta\) and \(\sin\theta\), reaching a ratio of trigonometric functions of \(\frac{\theta}{2}\).} \\
& {\(z = -i\cot\left(\frac{(2k+1)\pi}{22}\right)\)} & A1 & {cso, with \(\frac{1}{i} = -i\); any equivalent set of eleven values.} \\
& {\(z\) real \(\Rightarrow w\) real \(\Rightarrow w = -1 \Rightarrow z = 0\)\\[4pt] One real root (\(m = 11\)), so ten non-real} & E1 & {Identifies \(z = 0\) as the only real root, by \(w\) real or by \(\cot\left(\frac{m\pi}{22}\right) \neq 0\) for \(m \neq 11\), and concludes ten non-real roots.} \\
& {\(-i\cot\left(\pi - \frac{r\pi}{22}\right) = i\cot\left(\frac{r\pi}{22}\right)\)\\[4pt] \(\phantom{-i\cot\left(\pi - \frac{r\pi}{22}\right)} = -i\cot\left(-\frac{r\pi}{22}\right)\)\\[4pt] \(m = \pm1, \pm3, \pm5, \pm7, \pm9\)} & A1~AG & {Maps \(m = 13, \ldots, 21\) to \(-9, \ldots, -1\) by \(\cot(\pi - x) = -\cot x\) with cot odd, or by period \(\pi\), and states the printed set.} \\
\nopagebreak & & [6] & \\
13(iii) & {\((-i \times i)\cot^2\left(\frac{m\pi}{22}\right) = -i^2\cot^2\left(\frac{m\pi}{22}\right)\)\\[4pt] \(\phantom{(-i \times i)\cot^2\left(\frac{m\pi}{22}\right)} = \cot^2\left(\frac{m\pi}{22}\right)\)} & B1 & {Needs \(i^2 = -1\) used explicitly.} \\
\nopagebreak & & [1] & \\
13(iv) & {\(\cot^2\left(\frac{\pi}{22}\right)\cot^2\left(\frac{3\pi}{22}\right)\cdots\cot^2\left(\frac{9\pi}{22}\right) = 11\)} & M1 & {Pairs the roots for \(m\) and \(-m\) and uses (iii) to equate the product of the five \(\cot^2\) values to 11.} \\
& {\(\frac{1}{11}\)} & A1 & {\(\frac{1}{11}\) or awrt 0.0909.} \\
\nopagebreak & & [2] & \\
\nopagebreak \SetCell[c=2]{r} Total for Question 13 & & 11 & \\
\end{msscheme}

\end{document}
